\documentclass[11pt]{article}
\usepackage[margin=1in]{geometry}
\usepackage{amsmath,amssymb}
\usepackage{booktabs}
\usepackage{graphicx}
\usepackage{hyperref}
\usepackage{xcolor}
\hypersetup{colorlinks=true,linkcolor=blue,citecolor=blue,urlcolor=blue}
\title{\textbf{Three vertex models for epithelial morphogenesis in Plexus2}\\[2pt]
\large SPV, mid-surface AVM, and the monolayer AVM behind \texttt{mono\_V4\_tube\_flat}}
\author{Plexus2 prototype note --- \texttt{prototype/Tyssue}}
\date{2026-07-23}
\begin{document}
\maketitle

\begin{abstract}
This note fixes, in one place, the mathematics of the three cell-mechanics models the Plexus2 tissue prototype
uses, so that the \texttt{mono\_V4\_tube\_flat} simulation --- a flat epithelium that grows a tube by localized
proliferation --- can be studied on a sound footing. We give the energy, the equations of motion, the
topological update, and the Plexus2 operator for each: (A) the \emph{Self-Propelled Voronoi} (SPV) model
(cells are points), (B) the \emph{mid-surface Active Vertex Model} (AVM) used for the vesicle tubing runs
(cells are surface polygons with a lumen-wedge volume), and (C) the \emph{monolayer AVM} (cells are prisms
with their own apical/basal surfaces and a true 3D volume), which \texttt{mono\_V4} is built on. We then state
precisely how the three differ, and set up the study of \texttt{mono\_V4}.
\end{abstract}

\begin{figure}[h]
\centering
\includegraphics[width=\textwidth]{model_schematic.png}
\caption{The three cell-mechanics models. \textbf{(A)} Self-Propelled Voronoi: a cell is a \emph{point}
$r_i$; its polygon is the derived Voronoi cell; T1 is implicit (retessellate); motion adds self-propulsion
$v_0\hat n$. \textbf{(B)} Mid-surface AVM: the mesh \emph{vertices} $x_i$ are the DOF, cells are faces sharing
junctions; the per-cell volume is a lumen \emph{wedge} $v_f=\frac13 c_f\!\cdot\!N_f$ from the shell centre
$O$; T1 is an explicit flip. \textbf{(C)} Monolayer AVM (cross-section): each cell is a \emph{prism} with an
apical (outer) and basal (inner) surface offset $\pm h/2$ along the vertex normal, a true volume $A^{mid}h$
and surface $A^{ap}{+}A^{ba}{+}A^{lat}$; on a curved sheet the apical arc exceeds the basal arc, so surface
tension penalises curvature --- bending is \emph{emergent}.}
\label{fig:models}
\end{figure}

\section{Common ground}
All three are \emph{overdamped} (inertia-free) cell-mechanics models: a mechanical energy $U$ assigns a
force to each degree of freedom, and the tissue moves down that force. They differ in (i) \emph{what the
degrees of freedom are}, (ii) \emph{how a cell's shape/volume is measured}, and (iii) \emph{how neighbour
exchange (T1) happens}. The overdamped law is, generically,
\begin{equation}
\eta\,\frac{d\mathbf{q}}{dt} \;=\; -\,\frac{\partial U}{\partial \mathbf{q}} \;+\; \mathbf{f}^{\text{active}},
\label{eq:overdamped}
\end{equation}
with $\mathbf{q}$ the DOF, $\eta$ a friction, and $\mathbf{f}^{\text{active}}$ any non-conservative drive
(self-propulsion, extrusion). Plexus2 integrates this with a \emph{bounded} explicit Euler step (each
displacement capped at a fraction of the mean edge length) so the rollout stays stable \emph{and}
differentiable.

\section{Model A --- Self-Propelled Voronoi (SPV)}
\textbf{Degrees of freedom.} Each cell $i$ is a single \emph{point} $\mathbf{r}_i$. The tissue is the
\emph{Voronoi tessellation} of the points; cell $i$'s polygon (area $A_i$, perimeter $P_i$) is a derived
quantity, recomputed every step.

\textbf{Energy.} The standard 2D vertex-model shape energy evaluated on the Voronoi cells:
\begin{equation}
U_{\text{SPV}}=\sum_i \Big[\,K_A\,(A_i-A_0)^2 \;+\; K_P\,(P_i-P_0)^2\,\Big],
\qquad p_0 \equiv \frac{P_0}{\sqrt{A_0}}\ \text{(target shape index)}.
\label{eq:spv}
\end{equation}
The dimensionless target shape index $p_0$ sets a rigidity (jamming) transition: the tissue is a \emph{solid}
(finite energy barrier to rearrangement) for $p_0\lesssim 3.81$ and a \emph{fluid} above it
\cite{bi2015,bi2016}.

\textbf{Dynamics.} Overdamped motion plus a persistent self-propulsion of fixed speed $v_0$ along a heading
$\theta_i$ that rotationally diffuses (coefficient $D_r$):
\begin{equation}
\frac{d\mathbf{r}_i}{dt}=\mu\Big(-\nabla_i U_{\text{SPV}}\Big)+v_0\,\hat{\mathbf{n}}(\theta_i),
\qquad \frac{d\theta_i}{dt}=\sqrt{2D_r}\,\xi_i(t).
\label{eq:spvdyn}
\end{equation}

\textbf{T1 is \emph{implicit}.} Because the polygons are re-derived from the points each step, a neighbour
exchange happens automatically the instant two points' Voronoi cells stop sharing an edge --- there is no
explicit surgery. This is the model's great convenience and its central limitation: the topology is slaved to
the point positions, so it cannot represent a configuration the Voronoi rule forbids (e.g.\ a high-aspect
junctional tube, a lumen, a fourfold vertex held open).

\textbf{Plexus2 implementation.} \texttt{prototype/Turing\_vertex}: operator \texttt{voronoi\_tension}
(2D)/\texttt{vertex\_tension} on a \texttt{cell} set of points; the Voronoi is rebuilt each tick; parameters
\texttt{p0, K\_A, K\_P, mu, v0, Dr}. Reaction--diffusion couples on the Delaunay (neighbour) graph.

\section{Model B --- mid-surface Active Vertex Model (the tubing runs)}
\textbf{Degrees of freedom.} The \emph{vertices} $\mathbf{x}_i\in\mathbb{R}^3$ of an explicit half-edge mesh.
A cell is a \emph{face} (an ordered ring of vertices); neighbouring faces \emph{share} vertices and edges.
This is the true vertex model (AVM): the junctional network itself is the state, not a derived tessellation.

\textbf{Geometry.} Per face $f$: area vector by the Newell formula
$\mathbf{N}_f=\tfrac12\sum_{e\in f}\mathbf{x}_{s(e)}\times\mathbf{x}_{t(e)}$, area $A_f=\|\mathbf{N}_f\|$,
perimeter $P_f=\sum_{e\in f}\ell_e$, and a \emph{lumen-wedge volume} --- the signed volume of the pyramid
from the shell centre to the face,
\begin{equation}
v_f=\tfrac13\,\mathbf{c}_f\!\cdot\!\mathbf{N}_f,\qquad \mathbf{c}_f=\text{face centroid}.
\end{equation}
The wedge volume is a \emph{proxy} for the cell's contribution to the enclosed lumen; each cell carries its
own, so the term below resists local buckling (a distributed constraint, not one global lumen scalar).

\textbf{Energy.}
\begin{equation}
U_{\text{AVM}}=\sum_f\Big[K_A(A_f-A_f^0)^2+K_P(P_f-P_f^0)^2+\tfrac12\Gamma P_f^2+K_V(v_f-v_f^{\text{eq}})^2\Big]
\;+\;\Lambda\!\sum_e \ell_e \;+\;K_R\!\sum_i(\|\mathbf{x}_i\|-R_0)^2 .
\label{eq:avm}
\end{equation}
Terms: area elasticity, perimeter elasticity, cortical contractility $\tfrac12\Gamma P^2$, per-cell volume
elasticity $K_V$, line tension $\Lambda$ on edges, and an optional soft radial anchor $K_R$
\cite{farhadifar2007,okuda2013,okuda2015}.

\textbf{Dynamics \& topology.} Force $\mathbf{F}_i=-\partial U_{\text{AVM}}/\partial\mathbf{x}_i$ by one
autograd pass; bounded overdamped Euler, iterated \texttt{relax\_iters} times per frame toward force balance
(Okuda relaxes to a residual threshold $E_{\text{th}}$ instead --- see \S\ref{sec:mono}). T1 is
\emph{explicit}: \texttt{reconnect\_t1\_3d} flips any interior edge shorter than $\Delta\ell_{\text{th}}$ by a
local four-cell neighbour exchange, keeping the mesh closed (Euler $=2$). Division splits a face by an
edge-midpoint septum \emph{normal to the cell's longest axis}; the reversible network reconnection (RNR) of
\cite{okuda2013} is the 3D generalization.

\textbf{Plexus2 implementation.} \texttt{shape\_energy\_3d} (the ``default'' implementation of the
\texttt{shape\_energy\_3d} contract), \texttt{reconnect\_t1\_3d}, \texttt{divide\_3d}, on a \texttt{vertex}
set carrying the half-edge table. This is the model used for the vesicle tubing runs (\texttt{round\_*}).

\section{Model C --- monolayer AVM (behind \texttt{mono\_V4\_tube\_flat})}
\label{sec:mono}
\textbf{Motivation.} The mid-surface model represents a cell only by a polygon plus a wedge volume; it has no
\emph{thickness}, so it cannot express the apical/basal distinction, a true cell volume, or the bending
rigidity of a monolayer. Okuda's actual 3D model \cite{okuda2015,okuda2018} gives every cell its own volume
and surface. Model C is the minimal lift of the AVM to that: keep the mid-surface vertices as the only DOF,
but give each cell a \emph{thickness} and read off an apical and a basal surface.

\textbf{Geometry.} For thickness $h_j$ per cell and the area-weighted \emph{vertex normal}
$\mathbf{n}_i=\widehat{\sum_{f\ni i}\mathbf{N}_f}$, define the apical (outer) and basal (inner) vertices
\begin{equation}
\mathbf{a}_i=\mathbf{x}_i+\tfrac{H_i}{2}\,\mathbf{n}_i,\qquad
\mathbf{b}_i=\mathbf{x}_i-\tfrac{H_i}{2}\,\mathbf{n}_i,\qquad
H_i=\text{mean of }h_j\text{ over cells at }i .
\end{equation}
Each cell is then a \emph{prism} (apical cap $+$ basal cap $+$ lateral quads), with
\begin{equation}
v_j=A^{\text{mid}}_j\,h_j\quad(\text{cell volume}),\qquad
s_j=A^{\text{ap}}_j+A^{\text{ba}}_j+A^{\text{lat}}_j\quad(\text{cell surface}),
\end{equation}
where $A^{\text{ap}}_j,A^{\text{ba}}_j$ are the Newell areas of the offset caps and $A^{\text{lat}}_j$ sums
the lateral quad areas. \emph{Emergent bending}: because the caps are offset along the \emph{vertex} normal,
a curved sheet has $A^{\text{ap}}\neq A^{\text{ba}}$ (the convex side stretches), so surface tension
$\kappa_s s$ penalises curvature with \emph{no explicit bending term} --- the apical/basal area ratio matches
$((R{+}h/2)/(R{-}h/2))^2$ to $0.2\%$ in the validation.

\textbf{Energy (Okuda Eq.~3, plus a shape regularizer).}
\begin{equation}
U_{\text{mono}}=\sum_j\Big[\tfrac12 k_v(v_j-v_j^{\text{eq}})^2+\kappa_s\,s_j+\tfrac12\gamma P_j^2\Big]
\;+\;\Lambda\!\sum_e\ell_e\;+\;K_R\!\sum_i(\|\mathbf{x}_i\|-R_0)^2 .
\label{eq:mono}
\end{equation}
The first two terms are Okuda's per-cell volume elasticity and \emph{linear surface tension}. The
$\tfrac12\gamma P^2$ contractility is a \emph{cell-shape regularizer} standing in for the continuous
RNR/T1 remeshing Okuda relies on: without it the bare volume$+$surface energy shears or spikes under large
deformation.

\textbf{Per-cell stress (from the same energy).} The isotropic pressure and cortical tension read off Eq.~\ref{eq:mono} are
\begin{equation}
p_j=-\frac{\partial U}{\partial v_j}=2k_v\,(v_j^{\text{eq}}-v_j)\ \ (>0:\text{below-target volume, pushes out}),
\qquad
T_j=\frac{\partial U}{\partial P_j}=\gamma P_j+\Lambda .
\end{equation}
These are the fields \texttt{analyze\_forces.py} plots; the whole energy is torch-differentiable, so a loss
on stress can be optimised.

\textbf{\texttt{mono\_V4\_tube\_flat} specifics.} The substrate is a \emph{flat} open epithelium
(\texttt{build\_flat\_mesh}: a jittered planar Voronoi patch, boundary rim \emph{pinned}). A central
\emph{growth zone} that rides the advancing tip has its target volume $v^{\text{eq}}$ raised each round; the
tissue relaxes under Eq.~\ref{eq:mono} (bounded Euler, tangential Lloyd smoothing); oversized cells
\emph{divide in-plane}, the septum normal to the longest axis (Hertwig). Localized in-plane proliferation on
a rim-pinned sheet has nowhere to go but \emph{buckle out of plane}, extending a tube whose wall is the new
cells. There is no reaction--diffusion here --- the growth zone is prescribed geometrically.

\textbf{Plexus2 implementation.} \texttt{shape\_energy\_3d} with \texttt{implementation=monolayer}
(class \texttt{MonolayerShapeEnergy3D}) --- the plexus2 ``same biology, different numerics'' pattern: the
monolayer is a second numerical implementation of the same shape-energy operator contract. Geometry:
\texttt{monolayer\_geometry\_3d}; flat substrate: \texttt{flat\_mesh.py}; driver: \texttt{demo\_monolayer.py}
(\texttt{demo\_tube}). Division reuses the mid-surface topology ops (\texttt{divide\_face\_3d}).

\section{How the three differ}
\begin{center}\small
\begin{tabular}{@{}llll@{}}
\toprule
 & \textbf{A. SPV} & \textbf{B. mid-surface AVM} & \textbf{C. monolayer AVM} \\
\midrule
DOF & cell centres $\mathbf{r}_i$ & mesh vertices $\mathbf{x}_i$ & mesh vertices $\mathbf{x}_i$ ($+h_j$) \\
cell shape & Voronoi polygon (derived) & face polygon (explicit) & prism: apical+basal+lateral \\
volume & --- (2D) & lumen wedge $\tfrac13\mathbf{c}\!\cdot\!\mathbf{N}$ & true $A^{\text{mid}}h$ \\
surface & perimeter only & area+perimeter & apical+basal+lateral \\
bending & --- & none (or explicit $K_{\text{bend}}$) & \emph{emergent} (cap offset) \\
T1 / neighbour exch. & implicit (retessellate) & explicit RNR/T1 flip & explicit (inherits B's) \\
energy & Eq.~\ref{eq:spv} & Eq.~\ref{eq:avm} & Eq.~\ref{eq:mono} (Okuda 3) \\
can hold a tube? & no (Voronoi-forbidden) & yes (junctional net) & yes, with a true wall \\
\bottomrule
\end{tabular}
\end{center}

\noindent\textbf{The essential distinctions.}
\emph{A vs B:} SPV's topology is a slave of the point positions (implicit T1), so it is fast and always valid
but cannot represent a high-aspect tube or a lumen; the AVM makes the junctional network the state, paying
for explicit T1 surgery but able to hold a tube. \emph{B vs C:} the mid-surface AVM has no thickness --- its
``volume'' is a wedge from the shell centre and it has no apical/basal surfaces or bending; the monolayer AVM
gives each cell a real volume $A\,h$ and a real surface, so localized volume growth \emph{buckles} into a
tube as an \emph{equilibrium} shape (Okuda's growth-driven, quasi-static route), whereas the mid-surface tube
in the \texttt{round\_*} runs is held out by an added extrusion force (a forced, higher-stress protrusion).

\section{The \texttt{mono\_V4\_tube\_flat} study (what we will make sound)}
\texttt{mono\_V4} currently extends a lumpy, chaotic tube: growth is confined to a moving tip cap, cells
divide when oversized, and the sheet buckles --- but the wall is irregular and the dynamics are not
understood. The plan is to make this one simulation rigorous:
\begin{enumerate}\itemsep2pt
\item \textbf{Start from a working point.} \texttt{mono\_V3b\_buckle\_flat} (a clean localized buckle, modest
amplitude) is a controlled, understood equilibrium; growing the tube \emph{from} that buckle --- rather than
the current fast, uncontrolled extension --- is the sound way in.
\item \textbf{Levers.} $\kappa_s/k_v$ (surface tension vs volume stiffness $\to$ buckle amplitude and wall
tension), $h_0$ (thickness $\to$ emergent bending, tube diameter), $\gamma$ (cell rounding), the growth rate
and division threshold (quasi-static vs explosive), and the tip-cap size (tube diameter).
\item \textbf{Analysis.} Run \texttt{analyze\_forces.py} throughout: per-cell force $\|-\nabla U\|$, pressure
$p_j$, cortical tension $T_j$, cell migration and division --- to read the wall stress and check the
growth-driven regime is quasi-static (bounded, tube$\leftrightarrow$body-shared stress) rather than
explosive.
\end{enumerate}

\begin{thebibliography}{9}
\bibitem{bi2015} D.~Bi, J.~H.~Lopez, J.~M.~Schwarz, M.~L.~Manning (2015). \emph{A density-independent
rigidity transition in biological tissues.} Nat.\ Phys.\ 11, 1074.
\bibitem{bi2016} D.~Bi, X.~Yang, M.~C.~Marchetti, M.~L.~Manning (2016). \emph{Motility-driven glass and
jamming transitions in biological tissues.} Phys.\ Rev.\ X 6, 021011.
\bibitem{farhadifar2007} R.~Farhadifar, J.-C.~R\"oper, B.~Aigouy, S.~Eaton, F.~J\"ulicher (2007). \emph{The
influence of cell mechanics, cell-cell interactions, and proliferation on epithelial packing.} Curr.\ Biol.\
17, 2095.
\bibitem{okuda2013} S.~Okuda, Y.~Inoue, M.~Eiraku, Y.~Sasai, T.~Adachi (2013). \emph{Reversible network
reconnection model for simulating large deformation in dynamic tissue morphogenesis.} Biomech.\ Model.\
Mechanobiol.\ 12, 627.
\bibitem{okuda2015} S.~Okuda, Y.~Inoue, M.~Eiraku, T.~Adachi, Y.~Sasai (2015). \emph{Vertex dynamics
simulations of viscosity-dependent deformation during tissue morphogenesis.} Biomech.\ Model.\ Mechanobiol.\
14, 413.
\bibitem{okuda2018} S.~Okuda, T.~Miura, Y.~Inoue, T.~Adachi, M.~Eiraku (2018). \emph{Combining Turing and 3D
vertex models reproduces autonomous multicellular morphogenesis of the tissue.} Sci.\ Rep.\ 8, 2386.
\end{thebibliography}
\end{document}
